\documentclass[11pt,a4paper]{article}
\usepackage[utf8]{inputenc}
\usepackage[T1]{fontenc}
\usepackage[english]{babel}
\usepackage{amsmath,amssymb,amsthm}
\usepackage{geometry}
\usepackage{enumitem}
\usepackage{fancyhdr}
\usepackage{array}
\usepackage{longtable}
\usepackage{hyperref}

\geometry{margin=1in}

\providecommand{\mid}{\,|\,}
\providecommand{\Disct}{\operatorname{Disct.}}

\theoremstyle{plain}
\newtheorem{theorem}{Theorem}

\pagestyle{fancy}
\fancyhf{}
\fancyhead[L]{\small Cayley, \textit{Coll.\ Math.\ Papers}, Vol.\ V}
\fancyhead[R]{\small pp.\ 126--150}
\fancyfoot[C]{\thepage}

\title{\textbf{The Collected Mathematical Papers of Arthur Cayley}\\[6pt]
\Large Volume V (1861--1863)\\[12pt]
\large Pages 126--150}
\author{Arthur Cayley}
\date{}

\usepackage{graphicx}
\emergencystretch=3em
\tolerance=600
\begin{document}
\maketitle
\thispagestyle{empty}
\newpage
\setcounter{page}{126}

% =====================================================================
% Paper 330 (continued from p. 125) -- pp. 126--130
% =====================================================================
\section*{330. On Differential Equations and Umbilici \textit{(continued)}}

\noindent[Continuation from p.~125, mid-derivation in section VI.]

\medskip

\noindent or
\[
(2m^2 - m - 1)\,A^2 + (2m^2 - 3m + 1)\,y^2 - (2m^2 - 2m)\,A\,(x + 2\sqrt{\square}) = 0,
\]
and therefore
\[
A^2 + (2m - 1)\,y^2 + (2m^2 - 2m)\,A\,(x + 2\sqrt{\square}) = (2m^2 - m)(A^2 + y^2).
\]
Hence the term in $\{\ \}$ is
\[
= (2m^2 - m)(A^2 + y^2)(A + yp);
\]
or, what is the same thing, it is $= (4m^2 - 2m)\,A\sqrt{\square}\,(A + yp)$. Hence, restoring for $A^2 + (2m - 1)\,y^2$ its value $2mB$, we find
\[
\Theta = \frac{(2m - 1)\,B\sqrt{\square}}{A}\,(A + yp),
\]
\[
\frac{U'}{U} = \frac{2m - 1}{B}\,(A + yp).
\]

But, writing $U_1$, $U_2$ to denote the values corresponding to $+\sqrt{\square}$, $-\sqrt{\square}$ respectively, we have
\[
U_1' = \frac{(2m - 1)\,U_1}{B_1}\,(mx + yp + \sqrt{\square}),
\]
\[
U_2' = \frac{(2m - 1)\,U_2}{B_2}\,(mx + yp - \sqrt{\square});
\]
and thence
\[
U_1'\,U_2' = \frac{(2m - 1)^2\,U_1 U_2}{B_1 B_2}\,\bigl\{(mx + yp)^2 - \square\bigr\}.
\]

But we have
\[
U' = P' + Q'\sqrt{\square} + \frac{Q\square'}{2\sqrt{\square}} = \frac{1}{2\sqrt{\square}}\,(2Q'\square + Q\square' + 2P'\sqrt{\square}\,),
\]
and thence
\[
U_1'\,U_2' = -\frac{1}{4\square}\,\bigl\{(2Q'\square + Q\square')^2 - 4P'^2\square\bigr\},
\]
and moreover
\[
U_1\,U_2 = P^2 - Q^2\square = A_1 A_2\,(B_1 B_2)^{m-1},
\]
where
\[
A_1 A_2 = m^2 x^2 - \square = -y^2,
\]
\[
B_1 B_2 = (m x^2 + y^2)^2 - x^2 \square = y^2\{y^2 + (2m - 1)\,x^2\};
\]
and we thence find
\[
-\tfrac{1}{4}\square\,\bigl\{(2Q'\square + Q\square')^2 - 4P'^2\square\bigr\}
 = (2m - 1)^2\,A_1 A_2\,(B_1 B_2)^{m-1}\{y(p^2 - 1) + 2mxp\}
\]
\[
= -(2m - 1)^2\,y^{2m}\bigl[y^2 + (2m - 1)\,x^2\bigr]^{m-1}\{y(p^2 - 1) + 2mxp\}.
\]

\newpage
Hence, the derived equation being
\[
\Omega = Q^2\bigl\{(2Q'\square + Q\square')^2 - 4P'^2\square\bigr\} = 0,
\]
the last preceding equation becomes
\[
Q^2\,\square\,y^{2m-1}\bigl[y^2 + (2m - 1)\,x^2\bigr]^{m-1}\{y(p^2 - 1) + 2mxp\} = 0.
\]
Here, besides the factor $Q^2$ corresponding to the nodal curve, and the factor $\square$ corresponding to the cuspidal curve, we have the factors $y^{2m-1}$ and $[y^2 + (2m - 1)\,x^2]^{m-1}$; and, rejecting all these, the differential equation in its reduced form is
\[
y(p^2 - 1) + 2mxp = 0;
\]
and the required verification is effected. The occurrence of
\[
Q^2\,\square\,y^{2m-1}\bigl[y^2 + (2m - 1)\,x^2\bigr]^{m-1}
\]
as a factor in the complete derived equation would give rise to some further investigations, but I will not now enter on them.

I remark however that if $m = 1$, viz.\ if the integral equation be $\text{const.} = x + \sqrt{x^2 + y^2}$, or say $z = x + \sqrt{x^2 + y^2}$, or, what is the same thing,
\[
z^2 - 2zx - y^2 = 0,
\]
then observing that $y^2 + (2m - 1)\,x^2$ is here $= x^2 + y^2$ which is $= \square$, so that
\[
\square\,[y^2 + (2m - 1)\,x^2]^{m-1} = \square\cdot\square^0 = 1,
\]
the differential equation in its complete form is
\[
y(p^2 y + 2px - y) = 0;
\]
so that we have here the factor $y$ which divides out. The last-mentioned result is most readily obtained directly from the equation
\[
\Omega = Q^2\bigl\{(2Q'\square + Q\square')^2 - 4P'^2\square\bigr\} = 0,
\]
which is the derived equation corresponding to the integral equation $z = P + Q\sqrt{\square}$. We in fact have $P = x$, $Q = 1$, $\square = x^2 + y^2$, and the derived equation thus is
\[
(x + yp)^2 - (x^2 + y^2) = 0,
\]
that is, $y(p^2 y + 2px - y) = 0$.

I mention also, in connexion with the foregoing investigation, the integral equation
\[
z = x + \sqrt{2x^2 - y^2},\quad \text{or}\quad z^2 - 2zx - x^2 + y^2 = 0,
\]
for which the derived equation in its complete form is
\[
(2x - yp)^2 - (2x^2 - y^2) = 0,
\]
or, what is the same thing, $y^2 p^2 - 4xyp + 2x^2 + y^2 = 0$, and for which therefore there is no factor to divide out.

\bigskip
\begin{center}\textbf{VII.}\end{center}

The conics confocal with a given conic form a system similar in its properties to that of the curves of curvature of a quadric surface; and the theory of the last-mentioned system may be studied by means of the system of confocal conics. Consider then the equation
\[
\frac{x^2}{a^2 + z} + \frac{y^2}{b^2 + z} = 1,
\]
which, if $z$ be an arbitrary parameter, belongs to the conics confocal with the ellipse
$\frac{x^2}{a^2} + \frac{y^2}{b^2} = 1$. Treating $z$ as a coordinate, the equation represents a surface of the third order, which is such that its section by any plane parallel to the plane of $xy$ is a conic; and the confocal conics are the projections on the plane of $xy$, by lines parallel to the axis of $z$, of the sections of the surface.

The sections by the planes of $zx$, $zy$ are the parabolas $x^2 = z + a^2$ and $y^2 = z + b^2$ respectively. When $z > -b^2$, the ordinates in each parabola are real, and these ordinates give the semiaxes of the elliptic section. When $z > -a^2$, $< -b^2$, then only the parabola section in the plane of $zx$ has a real ordinate, and the sections are hyperbolic; and when $z < -a^2$, the section is altogether imaginary. The section in the plane $z = -b^2$ is the pair of coincident lines $y^2 = 0$, $z = -b^2$, and the section in the plane $z = -a^2$ is the pair of coincident lines $z = -a^2$, $x^2 = 0$; or, in other words, the plane $z + b^2 = 0$ touches the surface along the line $y = 0$, and the plane $z + a^2 = 0$ touches the surface along the line $x = 0$: this at once appears from the integral form
\[
(z + a^2)(z + b^2) - x^2(z + b^2) - y^2(z + a^2) = 0.
\]

The points $(z = -b^2,\ y = 0,\ x = \pm\sqrt{a^2 - b^2})$ and $(z = -a^2,\ x = 0,\ y = \pm\sqrt{b^2 - a^2})$ are conical points; the last two are however imaginary points on the surface. To find the nature of the surface about one of the first-mentioned two points, say the point $(z = -b^2,\ y = 0,\ x = \sqrt{a^2 - b^2})$, taking this point for the origin and writing therefore $\sqrt{a^2 - b^2} + x$, $y$ and $-b^2 + z$ in the place of $x$, $y$, $z$ respectively, the equation becomes
\[
(a^2 - b^2 + z)\,z - \bigl\{(a^2 - b^2) + 2x\sqrt{a^2 - b^2} + x^2\bigr\}\,z - (a^2 - b^2 + z)\,y^2 = 0,
\]
that is
\[
z^2 - 2zx\sqrt{a^2 - b^2} - (a^2 - b^2)\,y^2 - z(x^2 + y^2) = 0;
\]
so that there is a tangent cone the equation whereof is
\[
z^2 - 2zx\sqrt{a^2 - b^2} - (a^2 - b^2)\,y^2 = 0,
\]
or, as it may be written,
\[
(z - x\sqrt{a^2 - b^2})^2 - (a^2 - b^2)(x^2 + y^2) = 0.
\]

\newpage
The equation is that of a cone of the second order, meeting the plane of $zx$ in the lines $z = 0$, $z = 2x\sqrt{a^2 - b^2}$ (and therefore such that its sections parallel to the plane of $xy$ are parabolas), and meeting the plane of $yz$ in the lines $z = \pm y\sqrt{a^2 - b^2}$ (the origin being at the vertex of the cone or conical point of the surface).

Returning to the original origin, and to the equation of the surface written in the form
\[
z^2 + z(a^2 + b^2 - x^2 - y^2) + a^2 b^2 - b^2 x^2 - a^2 y^2 = 0,
\]
calling this for a moment $z^2 + 2Bz + C = 0$, the differential equation is $C'^2 - 4BB'C + 4CB'^2 = 0$; or, substituting, this is
\[
(b^2 x + a^2 yp)^2 - (a^2 + b^2 - x^2 - y^2)(x + yp)(b^2 x + a^2 yp) + (a^2 b^2 - b^2 x^2 - a^2 y^2)(x + yp)^2 = 0;
\]
or, reducing, this is
\[
(a^2 - b^2)\,xy\,\bigl\{xy(p^2 - 1) - (a^2 - b^2 - x^2 + y^2)\,p\bigr\} = 0,
\]
or say
\[
xy\,\bigl\{xy(p^2 - 1) - (a^2 - b^2 - x^2 + y^2)\,p\bigr\} = 0,
\]
where the factor $xy$ arises from the level lines $(z + b^2 = 0,\ y = 0)$ and $(z + a^2 = 0,\ x = 0)$. Throwing out this factor, the equation becomes
\[
xy(p^2 - 1) - (a^2 - b^2 - x^2 + y^2)\,p = 0,
\]
which is satisfied identically by $z + b^2 = 0$, $y = 0$, $x^2 = a^2 - b^2$. The first derived equation is
\[
(xp + y)(p^2 - 1) + 2(x - yp)\,p = 0,
\]
which for the values in question gives
\[
p(p^2 + 1) = 0,
\]
where the factor $p = 0$ corresponds to the section $y = 0$ by the plane $z + b^2 = 0$: and taking the conical point for origin, and observing that the polar of the line $x = 0$, $y = 0$ in regard to the tangent cone is $z - x\sqrt{a^2 - b^2} = 0$, then writing the equation of the tangent cone in the form
\[
(z - x\sqrt{a^2 - b^2})^2 - (a^2 - b^2)(x^2 + y^2) = 0,
\]
the two tangent planes through $(x = 0,\ y = 0)$ are given by the equation $x^2 + y^2 = 0$; and for these planes we have $p^2 + 1 = 0$. The factor $p^2 + 1 = 0$ determines therefore the directions of the envelope at the conical point.

\bigskip
\begin{center}\textbf{VIII.}\end{center}

In verification of the equation
\[
z = \tfrac{1}{2}k(x^2 + y^2) + \tfrac{1}{6}\,a x(x^2 + y^2)
\]
for a quadric surface in the neighbourhood of the umbilicus, I remark that, starting from the equation
\[
\frac{x^2}{a^2} + \frac{y^2}{b^2} + \frac{z^2}{c^2} = 1
\]
of an ellipsoid, and taking $\alpha,\ 0,\ \gamma$ as the coordinates of the umbilicus, and $\theta$ as the inclination to the axis of $x$ of the tangent to the principal section through the umbilicus, then transforming to the umbilicus as origin and the new axes through that point, viz.\ the axes of $x$, $z$ being the tangent and normal in the plane of $\alpha z$, and the axis of $y$ being at right angles to this (or in the direction of $b$), the equation becomes
\[
\frac{(\alpha + x\cos\theta - z\sin\theta)^2}{a^2} + \frac{y^2}{b^2} + \frac{(\gamma - x\sin\theta - z\cos\theta)^2}{c^2} = 1,
\]
or, expanding,
\[
\Bigl(\frac{\alpha^2}{a^2} + \frac{\gamma^2}{c^2} - 1\Bigr) + 2x\Bigl(\frac{\alpha\cos\theta}{a^2} - \frac{\gamma\sin\theta}{c^2}\Bigr) - 2z\Bigl(\frac{\alpha\sin\theta}{a^2} + \frac{\gamma\cos\theta}{c^2}\Bigr)
\]
\[
{}+ x^2\Bigl(\frac{\cos^2\theta}{a^2} + \frac{\sin^2\theta}{c^2}\Bigr) + \frac{y^2}{b^2} + z^2\Bigl(\frac{\sin^2\theta}{a^2} + \frac{\cos^2\theta}{c^2}\Bigr) - 2xz\sin\theta\cos\theta\Bigl(\frac{1}{a^2} - \frac{1}{c^2}\Bigr) = 0.
\]

But we have
\[
\alpha = a\,\sqrt{\frac{a^2 - b^2}{a^2 - c^2}},\quad \gamma = \frac{c\,\sqrt{b^2 - c^2}}{\sqrt{a^2 - c^2}},
\]
\[
\tan\theta = \frac{c\,\sqrt{a^2 - b^2}}{a\,\sqrt{b^2 - c^2}},
\]
and thence
\[
\sin\theta = \frac{c\,\sqrt{a^2 - b^2}}{b\,\sqrt{a^2 - c^2}},\quad \cos\theta = \frac{a\,\sqrt{b^2 - c^2}}{b\,\sqrt{a^2 - c^2}};
\]
and substituting these values, the equation becomes
\[
-\frac{2zb}{ca}\sqrt{a^2 - c^2} + \frac{x^2}{b^2} + \frac{y^2}{b^2} + z^2\,\frac{a^2 b^2 + c^2 b^2 - a^2 c^2}{a^2 b^2 c^2} - \frac{xz(a^2 - c^2)(b^2 - c^2 + b^2)}{a^2 c^2 b^2}\,\bigl(\,\ldots\,\bigr) = 0,
\]
or, what is the same thing,
\[
z = \frac{ca}{b\sqrt{a^2 - c^2}}\cdot\tfrac{1}{2}(x^2 + y^2) + \tfrac{1}{2}\,\bigl\{\text{terms in } x^2 z,\ y^2 z,\ z^2,\ xz\bigr\} + \cdots,
\]
whence approximately
\[
z = \tfrac{1}{2}\,\frac{ca}{b\sqrt{a^2 - c^2}}\,(x^2 + y^2),
\]
and thence to the third order in $x$, $y$,
\[
z = \tfrac{1}{2}\,\frac{ca}{b\sqrt{a^2 - c^2}}\,(x^2 + y^2) + \tfrac{1}{6}\,a x(x^2 + y^2),
\]
which is of the form in question.

\medskip
\noindent\textit{5, Downing Terrace, Cambridge, November 2, 1863.}

% =====================================================================
% Paper 331 -- pp. 131--132
% =====================================================================
\newpage
\section*{331. Analytical Theorem Relating to the Four Conics Inscribed in the Same Conic and Passing through the Same Three Points}

\noindent\textit{[From the Philosophical Magazine, vol.\ xxvii.\ (1864), pp.\ 42, 43.]}

\medskip

Imagine the four conics determined, and, selecting at pleasure any three of them, let their chords of contact with the given conic be taken for the axes of coordinates, or lines $x = 0$, $y = 0$, $z = 0$; then, taking for the equation of the given conic
\[
U = (a,\,b,\,c,\,f,\,g,\,h \mid x,\,y,\,z)^2 = 0,
\]
the equations of the selected three conics must be of the form $U + l x^2 = 0$, $U + m y^2 = 0$, $U + n z^2 = 0$, where $l$, $m$, $n$ are to be determined in such manner that these conics may have three common points; the resulting values of $l$, $m$, $n$, and of the coordinates of the three common points, that is, the three given points, will of course be functions of the coefficients $(a,\,b,\,c,\,f,\,g,\,h)$; and the equation of the fourth conic will be of the form $U + (i x + j y + k z)^2 = 0$.

There is no difficulty in carrying out the investigation: it is found that the coordinates of the given points must be taken to be
\[
(-f,\ g,\ h);\quad (f,\ -g,\ h);\quad (f,\ g,\ -h)
\]
respectively, and that, writing as usual
\[
K = abc - af^2 - bg^2 - ch^2 + 2fgh,
\]
the equations of the four conics are
\[
U + (K - abc)\,\frac{x^2}{f^2} = 0,
\]
\[
U + (K - abc)\,\frac{y^2}{g^2} = 0,
\]
\[
U + (K - abc)\,\frac{z^2}{h^2} = 0,
\]
\[
U + (K - abc)\Bigl(\frac{x}{f} + \frac{y}{g} + \frac{z}{h}\Bigr)^{\!2} = 0.
\]

It is in fact easy to verify directly that each of these conics passes through the three given points; but the equations may also be exhibited in the form proper for putting this in evidence. Putting for shortness
\[
X = \frac{y}{g} + \frac{z}{h},\quad Y = \frac{z}{h} + \frac{x}{f},\quad Z = \frac{x}{f} + \frac{y}{g},
\]
the equations of the sides of the triangle formed by the given points are $X = 0$, $Y = 0$, $Z = 0$, and the foregoing equations of the four conics may be expressed in the form
\begin{align*}
(-bg^2 - ch^2 + 2fgh)\,YZ\ &+\ bg^2\,ZX\ +\ ch^2\,XY = 0,\\[2pt]
af^2\,YZ\ +\ (-ch^2 - af^2 + 2fgh)\,ZX\ &+\ ch^2\,XY = 0,\\[2pt]
af^2\,YZ\ +\ bg^2\,ZX\ &+\ (-af^2 - bg^2 + 2fgh)\,XY = 0,\\[2pt]
(-bg^2 - ch^2 + 2fgh)\,YZ\ +\ (-ch^2 - af^2 + 2fgh)\,ZX\ &+\ (-af^2 - bg^2 + 2fgh)\,XY = 0,
\end{align*}
which is the required form.

\medskip
\noindent\textit{Cambridge, November 28, 1863.}

% =====================================================================
% Paper 332 -- pp. 133--134
% =====================================================================
\newpage
\section*{332. Analytical Theorem Relating to the Sections of a Quadric Surface}

\noindent\textit{[From the Philosophical Magazine, vol.\ xxvii.\ (1864), pp.\ 43, 44.]}

\medskip

The four sections $x = 0$, $y = 0$, $z = 0$, $w = 0$ of the quadric surface
\[
a x^2 + b y^2 + 6 xy\sqrt{ab} - c z^2 - d w^2 = 0
\]
are each of them touched by each of the four sections
\[
x\sqrt{2a} + y\sqrt{2b} + z\sqrt{c} + w\sqrt{d} = 0;
\]
where it is to be noticed that the radicals $\sqrt{2a}$, $\sqrt{2b}$ are such that their product is $= 2\sqrt{ab}$ if $\sqrt{ab}$ be the radical contained in the equation of the surface. There is of course no loss of generality in attributing a definite sign to the radical $\sqrt{2a}$; but upon this being done, the sign of the radical $\sqrt{2b}$ is determined, whereas the signs of $\sqrt{c}$ and $\sqrt{d}$ are severally arbitrary. We may if we please write the equation of any one of the last-mentioned sections in the form
\[
x\sqrt{2a} + y\sqrt{2b} + z\sqrt{c} + w\sqrt{d} = 0,
\]
it being understood that the radicals $\sqrt{2a}$, $\sqrt{2b}$ have each a determinate sign, but that the signs of $\sqrt{c}$ and $\sqrt{d}$ are each of them arbitrary.

To prove the theorem in question, it is enough to show
\begin{enumerate}[label=(\arabic*)]
\item that the sections $x = 0$, $x\sqrt{2a} + y\sqrt{2b} + z\sqrt{c} + w\sqrt{d} = 0$;
\item that the sections $z = 0$, $x\sqrt{2a} + y\sqrt{2b} + w\sqrt{d} = 0$,
\end{enumerate}
touch each other.

\medskip

1. The sections $x = 0$, $x\sqrt{2a} + y\sqrt{2b} + z\sqrt{c} + w\sqrt{d} = 0$ of the quadric surface $a x^2 + b y^2 + 6 xy\sqrt{ab} - c z^2 - d w^2 = 0$ will touch each other if, combining together the equations
\[
x = 0,\quad y\sqrt{2b} + z\sqrt{c} + w\sqrt{d} = 0,\quad b y^2 - c z^2 - d w^2 = 0,
\]
these give a twofold value (pair of equal values) for the ratios $y : z : w$. We in fact have
\begin{align*}
b y^2 - c z^2 - d w^2 &= b y^2 - c z^2 - (y\sqrt{2b} + z\sqrt{c}\,)^2\\
 &= -b y^2 - 2 c z^2 - 2 y z\sqrt{2bc}\\
 &= -(y\sqrt{b} + z\sqrt{2c}\,)^2;
\end{align*}
and the right-hand side being a perfect square, the condition of contact is satisfied.

\medskip

2. In like manner we have the system
\[
z = 0,\quad x\sqrt{2a} + y\sqrt{2b} + w\sqrt{d} = 0,\quad a x^2 + b y^2 + 6 xy\sqrt{ab} - d w^2 = 0,
\]
which gives
\begin{align*}
a x^2 + b y^2 + 6 xy\sqrt{ab} - d w^2 &= a x^2 + b y^2 + 6 xy\sqrt{ab} - (x\sqrt{2a} + y\sqrt{2b}\,)^2\\
 &= -a x^2 - b y^2 + 2 xy\sqrt{ab}\\
 &= -(x\sqrt{a} - y\sqrt{b}\,)^2;
\end{align*}
and here also, the right-hand side being a perfect square, the condition of contact is satisfied.

\medskip
\noindent\textit{Cambridge, November 28, 1863.}

% =====================================================================
% Paper 333 -- pp. 135--137
% =====================================================================
\newpage
\section*{333. Note on the Nodal Curve of the Developable Derived from the Quartic Equation $(a,\,b,\,c,\,d,\,e\,\mid\,t,\,1)^4 = 0$}

\noindent\textit{[From the Philosophical Magazine, vol.\ xxvii.\ (1864), pp.\ 437--440.]}

\medskip

Considering the coefficients $(a,\,b,\,c,\,d,\,e)$ as linear functions of the coordinates $x$, $y$, $z$, $w$, then the equation
\[
\Disct\,(a,\,b,\,c,\,d,\,e\,\mid\,t,\,1)^4 = 0,
\]
or, as it may be written,
\[
(ae - 4bd + 3c^2)^3 - 27\,(ace + 2bcd - ad^2 - b^2 e - c^3)^2 = 0
\]
represents, as is known, a developable surface or ``torse,'' having for its edge of regression (or cuspidal curve) the sextic curve the equations whereof are
\begin{align*}
ae - 4bd + 3c^2 &= 0,\\
ace + 2bcd - ad^2 - b^2 e - c^3 &= 0;
\end{align*}
and for its nodal curve, a curve the equations whereof (equivalent to two independent relations between the coordinates) are
\[
\frac{ac - b^2}{a} = \frac{ad - bc}{2b} = \frac{ae + 2bd - 3c^2}{6c} = \frac{be - cd}{2d} = \frac{ce - d^2}{e},
\]
or, as these may also be written,
\begin{align*}
a^2 d - 3abc + 2b^3 &= 0,\\
a^2 e + 2abd - 9ac^2 + 6b^2 c &= 0,\\
abe - 3acd + 2b^2 d &= 0,\\
ad^2 - b^2 e &= 0,\\
ade - 3bce + 2bd^2 &= 0,\\
ae^2 + 2bde - 9c^2 e + 6cd^2 &= 0,\\
be^2 - 3cde + 2d^3 &= 0;
\end{align*}
which curve is in fact an excubo-quartic, --- viz.\ a quartic curve the partial intersection of a quadric surface and a cubic surface, having in common two non-intersecting right lines. To show that this is so, I remark that the coefficients $a$, $b$, $c$, $d$, $e$, qu\^a linear functions of the four coordinates, satisfy a linear equation which may be taken to be
\[
a + b + c + d + e = 0;
\]
this being so, the first form shows that the curve in question lies on the quadric surface
\[
ac - b^2 + \tfrac{1}{2}(ad - bc) + \tfrac{1}{6}(ae + 2bd - 3c^2) + \tfrac{1}{2}(be - cd) + ce - d^2 = 0,
\]
or, as this equation may also be written,
\[
c(a + b + c + d + e) - b^2 + \tfrac{1}{2}ad + \tfrac{1}{6}(ae + 2bd) + \tfrac{1}{2}be - d^2 = 0.
\]
Substituting for $c$ its value, this equation is
\[
-(a + e + b + d)(\tfrac{1}{2}a + \tfrac{1}{2}e) - b^2 + \tfrac{1}{2}ad + \tfrac{1}{6}(ae + 2bd) + \tfrac{1}{2}be - d^2 = 0,
\]
or, what is the same thing,
\[
9(a + e + b + d)(a + e) + 6(b^2 + d^2) - 3(ad + be) - (ae + 2bd) = 0.
\]
Hence, finally, the equation of the quadric surface is
\[
9a^2 + 17ae + 9e^2 + 6b^2 - 2bd + 6d^2 + 9ab + 9de + 6ad + 6be = 0;
\]
and the curve lies also on the cubic surface
\[
ad^2 - b^2 e = 0.
\]

It only remains to show that these surfaces have in common two right lines, and to find the equations of these lines.

The cubic surface is a skew surface or ``scroll'' such that the equations of any generating line are $d - \theta b = 0$, $e - \theta^2 a = 0$, where $\theta$ is an arbitrary parameter. But considering the two lines
\[
(d - \theta_1 b = 0,\ e - \theta_1^2 a = 0),\quad (d - \theta_2 b = 0,\ e - \theta_2^2 a = 0),
\]
the general equation of the quadric surface through these two lines may be written
\begin{align*}
&A\,(d - \theta_1 b)(d - \theta_2 b)\\
{}+{}&B\,(e - \theta_1^2 a)(e - \theta_2^2 a)\\
{}+{}&C\,\bigl\{(d - \theta_1 b)(e - \theta_2^2 a) + (d - \theta_2 b)(e - \theta_1^2 a)\bigr\}\\
{}+{}&\frac{D}{\theta_1 - \theta_2}\,\bigl\{(d - \theta_1 b)(e - \theta_2^2 a) - (d - \theta_2 b)(e - \theta_1^2 a)\bigr\} = 0,
\end{align*}
or, expanding and reducing,
\begin{align*}
&A\,\bigl\{d^2 - (\theta_1 + \theta_2)\,bd + \theta_1\theta_2\,b^2\bigr\}\\
{}+{}&B\,\bigl\{e^2 - (\theta_1^2 + \theta_2^2)\,ea + \theta_1^2\theta_2^2\,a^2\bigr\}\\
{}+{}&C\,\bigl\{2de - (\theta_1^2 + \theta_2^2)\,ad - (\theta_1 + \theta_2)\,be + \theta_1\theta_2(\theta_1 + \theta_2)\,ab\bigr\}\\
{}+{}&D\,\bigl\{(\theta_1 + \theta_2)\,ad - be - \theta_1\theta_2\,ab\bigr\} = 0,
\end{align*}
which, if $\theta_1$, $\theta_2$ are the roots of the equation $\theta^2 - \tfrac{1}{3}\theta + 1 = 0$, and therefore $\theta_1 + \theta_2 = \tfrac{1}{3}$, $\theta_1\theta_2 = 1$, and $\theta_1^2 + \theta_2^2 = -\tfrac{17}{9}$, is
\begin{align*}
&A\,\bigl(d^2 - \tfrac{1}{3}\,db + b^2\bigr)\\
{}+{}&B\,\bigl(e^2 + \tfrac{17}{9}\,ae + a^2\bigr)\\
{}+{}&C\,\bigl(2de + \tfrac{17}{9}\,ad - \tfrac{1}{3}\,be + \tfrac{1}{3}\,ab\bigr)\\
{}+{}&D\,\bigl(\tfrac{1}{3}\,ad - be - ab\bigr) = 0.
\end{align*}
Putting $A = 6$, $B = 9$, $C = \tfrac{3}{2}$, $D = -\tfrac{9}{2}$, this is
\begin{align*}
&9\,(a^2 + \tfrac{17}{9}\,ae + e^2)\\
{}+{}&6\,(b^2 - \tfrac{1}{3}\,bd + d^2)\\
{}+{}&\tfrac{3}{2}\,(\tfrac{1}{3}\,ab + 2de + \tfrac{17}{9}\,ad - \tfrac{1}{3}\,be)\\
{}+{}&\tfrac{9}{2}\,(ab - \tfrac{1}{3}\,ad + be) = 0,
\end{align*}
which is the before-mentioned quadric surface; hence the quadric surface and the cubic surface intersect in the two lines
\[
(d - \theta_1 b = 0,\ e - \theta_1^2 a = 0),\quad (d - \theta_2 b = 0,\ e - \theta_2^2 a = 0)
\]
(where $\theta_1$, $\theta_2$ are the roots of the quadratic equation $\theta^2 - \tfrac{1}{3}\theta + 1 = 0$); and they consequently intersect also in an excubo-quartic curve, which is the theorem required to be proved.

\medskip
\noindent\textit{Blackheath, March 26, 1864.}

% =====================================================================
% Paper 334 -- pp. 138--140
% =====================================================================
\newpage
\section*{334. Note on the Theory of Cubic Surfaces}

\noindent\textit{[From the Philosophical Magazine, vol.\ xxvii.\ (1864), pp.\ 493--496.]}

\medskip

The equation
\[
AX^3 + BY^3 + 6\,CRST = 0,
\]
where $X + Y + R + S + T = 0$, represents a cubic surface of a special form, viz.\ each of the planes $R = 0$, $S = 0$, $T = 0$ is a triple tangent plane meeting the surface in three lines which pass through a point\footnote{The tangent plane of a surface intersects the surface in a curve having at the point of contact a double point, and in like manner a triple tangent plane intersects the surface in a curve with three double points, viz.\ each point of contact is a double point; there is not in general any triple tangent plane such that the three points of contact come together, or (what is the same thing) there is not in general any tangent plane intersecting the surface in a curve having at the point of contact a triple point. A surface may, however, have the kind of singularity just referred to, viz.\ a tangent plane intersecting the surface in a curve having at the point of contact a triple point; such tangent plane may be termed a `tritom' tangent plane, and its point of contact a `tritom' point: for a cubic surface the intersection by a tritom tangent plane is of course a system of three lines meeting in the tritom point. The tritom singularity is sibi-reciprocal; it is, I think, a singularity which should be considered in the theory of reciprocal surfaces.}; and, moreover, the three planes $AX^3 + BY^3 = 0$ are triple tangent planes intersecting in a line. It is worth noticing that the equation of the surface may also be written
\[
a x^3 + b y^3 + c(u^3 + v^3 + w^3) = 0,
\]
where $x + y + u + v + w = 0$. In fact, the coordinates satisfying the foregoing linear equations respectively, we have to show that the equation
\[
AX^3 + BY^3 + 6\,CRST = a x^3 + b y^3 + c(u^3 + v^3 + w^3)
\]
may be identically satisfied. We have
\begin{align*}
a x^3 + b y^3 + c\,(u^3 + v^3 + w^3)
 &= a x^3 + b y^3 + c\bigl[(u + v + w)^3 - 3(v + w)(w + u)(u + v)\bigr]\\
 &= a x^3 + b y^3 - c(x + y)^3 - 3c(v + w)(w + u)(u + v),
\end{align*}
which is to be
\[
= AX^3 + BY^3 + 6\,CRST;
\]
and we may find $X$, $Y$, $R$, $S$, $T$, linear functions of $x$, $y$, $u$, $v$, $w$, so as to satisfy these equations, and so that in virtue of
\[
x + y + u + v + w = 0,
\]
we shall have also $X + Y + R + S + T = 0$. For, assuming
\begin{align*}
AX^3 + BY^3 &= a x^3 + b y^3 - c(x + y)^3,\\
X + Y &= x + y,\\
R &= \tfrac{1}{2}(v + w),\quad C = -4c,\\
S &= \tfrac{1}{2}(w + u),\\
T &= \tfrac{1}{2}(u + v),
\end{align*}
we have identically
\begin{align*}
AX^3 + BY^3 + 6\,CRST &= a x^3 + b y^3 - c(x + y)^3 - 3c(v + w)(w + u)(u + v),\\
X + Y + R + S + T &= x + y + u + v + w,
\end{align*}
and thus it only remains to show that we can find $X$, $Y$ linear functions of $x$, $y$, such that
\begin{align*}
AX^3 + BY^3 &= a x^3 + b y^3 - c(x + y)^3,\\
X + Y &= x + y.
\end{align*}
This is always possible; in fact if
\[
U = a x^3 + b y^3 - c(x + y)^3,
\]
then taking $\Phi$ for the cubicovariant, and $\square$ for the discriminant of $U$, we have $\tfrac{1}{2}(\Phi + \sqrt{\square}\,U)$, $\tfrac{1}{2}(\Phi - \sqrt{\square}\,U)$ each a perfect cube, say
\[
\tfrac{1}{2}(\Phi + \sqrt{\square}\,U) = (\lambda x + \mu y)^3,\quad \tfrac{1}{2}(\Phi - \sqrt{\square}\,U) = (\nu x + \rho y)^3;
\]
and we then have
\[
U = \tfrac{1}{\sqrt{\square}}\bigl[(\lambda x + \mu y)^3 - (\nu x + \rho y)^3\bigr] = AX^3 + BY^3,
\]
which is satisfied by
\[
X = l(\lambda x + \mu y),\quad Y = m(\nu x + \rho y),
\]
if
\[
A l^3 = \frac{1}{\sqrt{\square}},\quad B m^3 = -\frac{1}{\sqrt{\square}}.
\]

\newpage
The equation $X + Y = x + y$ then gives
\begin{align*}
l\lambda + m\nu &= 1,\\
l\mu + m\rho &= 1,
\end{align*}
which give the values of $l$ and $m$, and thence the values of $A$ and $B$; and collecting all the equations, we have
\begin{align*}
X &= \frac{\rho - \nu}{\lambda\rho - \mu\nu}\,(\lambda x + \mu y),\quad A = \frac{1}{\sqrt{\square}}\Bigl(\frac{\lambda\rho - \mu\nu}{\rho - \nu}\Bigr)^{\!3},\\[2pt]
Y &= -\frac{\mu - \lambda}{\lambda\rho - \mu\nu}\,(\nu x + \rho y),\quad B = \frac{1}{\sqrt{\square}}\Bigl(\frac{\lambda\rho - \mu\nu}{\mu - \lambda}\Bigr)^{\!3},\\[2pt]
R &= \tfrac{1}{2}(v + w),\quad C = -4c,\\
S &= \tfrac{1}{2}(w + u),\\
T &= \tfrac{1}{2}(u + v),
\end{align*}
where
\[
\lambda x + \mu y = \bigl[\tfrac{1}{2}(\Phi + \sqrt{\square}\,U)\bigr]^{\frac{1}{3}},\quad \nu x + \rho y = \bigl[\tfrac{1}{2}(\Phi - \sqrt{\square}\,U)\bigr]^{\frac{1}{3}},
\]
($\Phi$, $\square$ being respectively the cubicovariant and the discriminant of $U = a x^3 + b y^3 - c(x + y)^3$), for the formulae of the transformation
\begin{align*}
AX^3 + BY^3 + 6\,CRST &= a x^3 + b y^3 + c(u^3 + v^3 + w^3),\\
X + Y + R + S + T &= x + y + u + v + w.
\end{align*}

The equation $a x^3 + b y^3 + c(u^3 + v^3 + w^3) = 0$, where $x + y + u + v + w = 0$, presents over the other form the advantage that it is included as a particular case under the equation $a x^3 + b y^3 + c u^3 + d v^3 + e w^3 = 0$ (where $x + y + u + v + w = 0$) employed by Dr Salmon as the canonical form of equation for the general cubic surface.

\medskip
\noindent\textit{5, Downing Terrace, Cambridge, April 29, 1864.}

% =====================================================================
% Paper 335 -- pp. 141--150 (introduction + Table I)
% =====================================================================
\newpage
\section*{335. Tables des Formes Quadratiques Binaires pour les D\'eterminants N\'egatifs depuis $D = -1$ jusqu'\`a $D = -100$, pour les D\'eterminants Positifs non Carr\'es depuis $D = 2$ jusqu'\`a $D = 99$, et pour les Treize D\'eterminants N\'egatifs Irr\'eguliers qui se trouvent dans le Premier Millier}

\noindent\textit{[From the Journal f\"ur die reine und angewandte Mathematik (Crelle), tom.\ LX.\ (1862), pp.\ 357--372.]}

\medskip

Les tables suivantes sont arrang\'ees de la mani\`ere prescrite dans les ``\textit{Disquisitiones arithmeticae}.'' Dans le m\'emoire de Lejeune Dirichlet ``Recherches sur diverses applications de l'analyse \`a la th\'eorie des nombres,'' tom.\ XIX (1839), p.~338 de ce Journal, on trouve un tableau dans lequel les r\`egles qui servent \`a former les caract\`eres des genres sont r\'esum\'ees. Soit $D = PS^2$ ou $2PS^2$, $S^2$ d\'esignant le plus grand carr\'e que $D$ contient, et $P$ un nombre impair; soient de plus $p,\,p',\,p''\ldots$ les facteurs premiers in\'egaux de $P$ et $r,\,r',\,r''\ldots$ les nombres premiers impairs qui divisent $S$ sans diviser $P$; \'ecrivons enfin pour abr\'eger $\delta = (-1)^{\frac{m-1}{2}}$, $\epsilon = (-1)^{\frac{m^2-1}{8}}$. Cela pos\'e on trouve \`a l'endroit cit\'e le tableau suivant:

\medskip

\noindent\textit{Premier cas}, $D = PS^2,\ P \equiv 1 \pmod{4}$:
\[
\begin{array}{l|l|l}
S \equiv 1\ (\bmod\ 2) & \tfrac{m}{p},\ \tfrac{m}{p'},\ \ldots & \\
S \equiv 2\ (\bmod\ 4) & \tfrac{m}{p},\ \tfrac{m}{p'},\ \ldots & \epsilon,\ \tfrac{m}{r},\ \tfrac{m}{r'},\ \ldots\\
S \equiv 0\ (\bmod\ 4) & \tfrac{m}{p},\ \tfrac{m}{p'},\ \ldots & \delta,\ \epsilon,\ \tfrac{m}{r},\ \tfrac{m}{r'},\ \ldots
\end{array}
\]

\noindent\textit{Deuxi\`eme cas}, $D = PS^2,\ P \equiv 3 \pmod{4}$:
\[
\begin{array}{l|l|l}
S \equiv 1\ (\bmod\ 2) & \delta,\ \tfrac{m}{p},\ \tfrac{m}{p'},\ \ldots & \tfrac{m}{r},\ \tfrac{m}{r'},\ \ldots\\
S \equiv 2\ (\bmod\ 4) & \delta,\ \tfrac{m}{p},\ \tfrac{m}{p'},\ \ldots & \delta\epsilon,\ \tfrac{m}{r},\ \tfrac{m}{r'},\ \ldots\\
S \equiv 0\ (\bmod\ 4) & \delta,\ \tfrac{m}{p},\ \tfrac{m}{p'},\ \ldots & \epsilon,\ \tfrac{m}{r},\ \tfrac{m}{r'},\ \ldots
\end{array}
\]

\noindent\textit{Troisi\`eme cas}, $D = 2PS^2,\ P \equiv 1 \pmod{4}$:
\[
\begin{array}{l|l|l}
S \equiv 1\ (\bmod\ 2) & \tfrac{m}{p},\ \tfrac{m}{p'},\ \ldots & \epsilon\\
S \equiv 0\ (\bmod\ 2) & \tfrac{m}{p},\ \tfrac{m}{p'},\ \ldots & \delta\epsilon,\ \tfrac{m}{r},\ \tfrac{m}{r'},\ \ldots
\end{array}
\]

\noindent\textit{Quatri\`eme cas}, $D = 2PS^2,\ P \equiv 3 \pmod{4}$:
\[
\begin{array}{l|l|l}
S \equiv 1\ (\bmod\ 2) & \delta,\ \tfrac{m}{p},\ \tfrac{m}{p'},\ \ldots & \delta\epsilon\\
S \equiv 0\ (\bmod\ 2) & \delta,\ \tfrac{m}{p},\ \tfrac{m}{p'},\ \ldots & \epsilon,\ \tfrac{m}{r},\ \tfrac{m}{r'},\ \ldots
\end{array}
\]

Dans ce tableau la notation $\tfrac{m}{p}$, dans laquelle j'omets les parenth\`eses usit\'ees, signifie le caract\`ere d'un nombre quelconque $m$ par rapport au nombre premier impair $p$, c.-\`a-d.\ que $m$ est r\'esidu ou non r\'esidu de $p$ selon que $\tfrac{m}{p} = +1$ ou $= -1$; de m\^eme $\delta$ est le caract\`ere de $m$ par rapport au nombre 4, savoir $m \equiv 1$ ou $3 \pmod 4$ selon que $\delta = +1$ ou $= -1$; enfin $\epsilon$, $\delta\epsilon$ sont les caract\`eres de $m$ par rapport au nombre 8, savoir $m \equiv 1$ ou 7 $(\bmod\ 8)$ pour $\epsilon = +1$, $\equiv 3$ ou 5 $(\bmod\ 8)$ pour $\epsilon = -1$; $m \equiv 1$ ou 3 $(\bmod\ 8)$ pour $\delta\epsilon = +1$, $\equiv 5$ ou 7 $(\bmod\ 8)$ pour $\delta\epsilon = -1$. Si pour un d\'eterminant donn\'e on veut former au moyen de ce tableau les caract\`eres des genres, on prend la ligne horizontale qui convient \`a ce d\'eterminant; \`a tous les caract\`eres $\tfrac{m}{p},\ \tfrac{m}{p'},\ \ldots\ \tfrac{m}{r},\ \tfrac{m}{r'},\ \ldots,\ \delta,\ \epsilon,\ \delta\epsilon$ qui se trouvent dans la ligne horizontale, on attribue les signes $+$ ou $-$ \`a volont\'e, avec cette restriction cependant que le signe compos\'e des signes qui se trouvent dans la premi\`ere partie de la ligne dont il s'agit soit positif. Si par exemple le d\'eterminant donn\'e est $D = -35$, on a $D = -35 = PS^2$, $P = -35 \equiv 1 \pmod 4$, $S = 1 \equiv 1 \pmod 2$, les nombres $p,\,p',\,\ldots$ sont 5, 7, et les signes que l'on doit consid\'erer sont $\tfrac{m}{5},\ \tfrac{m}{7}$. De l\`a on obtient les caract\`eres
\[
\begin{array}{c|c|c|c}
\tfrac{m}{5} & \tfrac{m}{7} & & \\ \hline
+ & + & & \\
- & - & &
\end{array}
\]
il y a donc deux genres de l'ordre proprement primitif. Dans le cas dont il s'agit (et en g\'en\'eral pour $D \equiv 1 \pmod 4$) il y a un ordre improprement primitif avec des genres qui ont les caract\`eres identiques \`a ceux des genres de l'ordre proprement primitif, les caract\`eres se rapportant dans ce cas \`a la moiti\'e d'un nombre quelconque repr\'esent\'e par la forme.

Pour faciliter l'impression des tables j'ai introduit deux nouvelles lettres $\alpha$ et $\beta$ dont voici la d\'efinition. Pour tous les d\'eterminants auxquels se rapportent mes tables, c.-\`a-d.\ pour les d\'eterminants n\'egatifs quelconques et positifs non-carr\'es depuis $-100$ jusqu'\`a $+99$ ainsi que pour les d\'eterminants n\'egatifs irr\'eguliers du premier millier, le nombre des facteurs premiers d\'esign\'es ci-dessus par les lettres $p,\,p',\,p''\ldots,\,r,\,r',\,r''\ldots$ n'exc\`ede pas deux. Soit donc $q$ le plus petit de ces facteurs premiers et $q'$ le plus grand lorsqu'il y en a deux, je d\'esigne par $\alpha$ le caract\`ere $\tfrac{m}{q}$ et par $\beta$ le caract\`ere $\tfrac{m}{q'}$.

Dans la colonne relative \`a la composition et portant l'inscription $Cp$, je repr\'esente comme \`a l'ordinaire par l'unit\'e la forme principale, par la lettre $c$ une forme qui produit par la duplication la forme principale, par les lettres $d$, $e$, $\ldots$ des formes qui la produisent par la triplication, la quadruplication, etc., de mani\`ere que l'on ait $c^2 = 1$, $d^3 = 1$, $e^4 = 1$, $f^5 = 1$, $g^6 = 1$, $h^7 = 1$, $i^8 = 1$, $j^9 = 1$, etc. Les notations $d$, $d_1$ par exemple signifient deux formes diff\'erentes dont chacune produit par la triplication la forme principale. Je repr\'esente de plus par $a$ la forme principale de l'ordre improprement primitif et par $ac$, $ad$, $\ldots$ des formes qui produisent $a$ par la duplication, la triplication, etc. Dans l'\'enum\'eration des classes, j'ai toujours \'ecrit en premier lieu l'ordre proprement primitif, en le faisant suivre apr\`es un trait de s\'eparation par l'ordre improprement primitif lorsqu'il existe. Dans chacun des deux ordres les divers genres se trouvent s\'epar\'es les uns des autres par des traits subordonn\'es.

Pour les d\'eterminants positifs les p\'eriodes sont donn\'ees par une abr\'eviation facile \`a comprendre. Chaque forme de la p\'eriode ayant son dernier coefficient \'egal au premier de la suivante, cette valeur identique n'a \'et\'e imprim\'ee qu'une fois; de plus les coefficients ext\'erieurs $a$, $c$ des formes $(a,\,b,\,c)$ ont \'et\'e distingu\'es des coefficients $b$ en imprimant ces derniers en caract\`eres plus petits. Ainsi pour le d\'eterminant 7 la p\'eriode de la classe principale $(1,\,0,\,-7)$ est donn\'ee par les nombres
\[
1,\ 2,\ -3,\ 1,\ 2,\ 1,\ -3,\ 2,\ 1
\]
qui repr\'esentent la s\'erie des formes
\[
(1,\,2,\,-3),\ (-3,\,1,\,2),\ (2,\,1,\,-3),\ (-3,\,2,\,1).
\]

\medskip
\noindent\textit{Londres, 6 Novembre, 1860.}

\bigskip

\subsection*{Table I. Formes quadratiques binaires ayant pour d\'eterminants les nombres n\'egatifs depuis $D = -1$ jusqu'\`a $D = -100$.}

\medskip

\noindent\textit{[Each entry below lists, for the given negative determinant $D$: the reduced forms (one per class), the genre-characters in the columns $\alpha,\ \beta,\ \delta,\ \epsilon,\ \delta\epsilon$ as required by the case, and the composition symbol $Cp$. The transcription preserves Cayley's reduced-form coefficients $(a,\,b,\,c)$; characters and composition symbols are abbreviated to the table headings used throughout.]}

\begin{longtable}{r|p{0.55\textwidth}|l|l}
$D$ & Reduced forms (one per class) & Characters & $Cp$ \\
\hline
\endhead
$-1$ & $(1,0,1)$ & $+$ & $1$ \\
$-2$ & $(1,0,2)$ & $+$ & $1$ \\
$-3$ & $(1,0,3)$ & $+$ & $1$ \\
$-4$ & $(1,0,4);\ (2,1,2)$ & $+$ & $1,\ c$ \\
$-5$ & $(1,0,5);\ (2,1,3)$ & $++$ & $1,\ c$ \\
$-6$ & $(1,0,6);\ (2,0,3)$ & $++$ & $1,\ c$ \\
$-7$ & $(1,0,7)$ & $++$ & $1$ \\
$-8$ & $(1,0,8);\ (3,1,3)$ & $++$ & $1,\ c$ \\
$-9$ & $(1,0,9);\ (2,1,5)$ & $++$ & $1,\ c$ \\
$-10$ & $(1,0,10);\ (2,0,5)$ & $++$ & $1,\ c$ \\
$-11$ & $(1,0,11);\ (3,1,4);\ (3,-1,4)$ & $+$ & $1,\ d,\ d^{2}$ \\
$-12$ & $(1,0,12);\ (3,0,4);\ (2,1,6)$ & $++$ & $1,\ c,\ c$ \\
$-13$ & $(1,0,13);\ (2,1,7)$ & $++$ & $1,\ c$ \\
$-14$ & $(1,0,14);\ (2,0,7);\ (3,1,5);\ (3,-1,5)$ & $++$ & $1,\ c,\ e,\ e^{3}$ \\
$-15$ & $(1,0,15);\ (3,0,5);\ (2,1,8)$ & $+++$ & $1,\ c,\ ac$ \\
$-16$ & $(1,0,16);\ (4,1,4);\ (4,-1,4)$ & $++$ & $1,\ c,\ c$ \\
$-17$ & $(1,0,17);\ (2,1,9);\ (3,1,6);\ (3,-1,6)$ & $++$ & $1,\ c,\ e,\ e^{3}$ \\
$-18$ & $(1,0,18);\ (2,0,9);\ (3,1,6\,\text{var.})$ & $++$ & $1,\ c$ \\
$-19$ & $(1,0,19);\ (4,1,5);\ (4,-1,5)$ & $+$ & $1,\ d,\ d^{2}$ \\
$-20$ & $(1,0,20);\ (4,2,6);\ (2,0,10);\ (5,0,4)$ & $++$ & $1,\ c,\ c$ \\
$-21$ & $(1,0,21);\ (2,1,11);\ (3,0,7);\ (5,2,5);\ (3,-1,7),\ (5,-2,5)$ & $++$ & $1,\ c,\ \ldots$ \\
$-22$ & $(1,0,22);\ (2,0,11)$ & $++$ & $1,\ c$ \\
$-23$ & $(1,0,23);\ (2,1,12);\ (3,1,8);\ (3,-1,8);\ (4,1,6);\ (4,-1,6)$ & $+$ & $1,\ d,\ d^{2},\ d_1,\ \ldots$ \\
$-24$ & $(1,0,24);\ (2,0,12);\ (3,0,8);\ (4,1,7);\ (4,-1,7)$ & $++$ & $1,\ c,\ c,\ \ldots$ \\
$-25$ & $(1,0,25);\ (5,0,5)$ & $++$ & $1,\ c$ \\
$-26$ & $(1,0,26);\ (2,0,13);\ (3,1,9);\ (3,-1,9);\ (5,2,6);\ (5,-2,6)$ & $++$ & $1,\ c,\ \ldots$ \\
$-27$ & $(1,0,27);\ (3,0,9);\ (4,1,7);\ (4,-1,7)$ & $+$ & $1,\ d,\ d^{2}$ \\
$-28$ & $(1,0,28);\ (4,0,7);\ (2,1,15);\ (5,2,6);\ (5,-2,6)$ & $++$ & $1,\ c,\ \ldots$ \\
$-29$ & $(1,0,29);\ (2,1,15);\ (5,1,6);\ (5,-1,6);\ (3,1,10);\ (3,-1,10);\ \ldots$ & $++$ & $1,\ d,\ \ldots$ \\
$-30$ & $(1,0,30);\ (2,0,15);\ (3,0,10);\ (5,0,6)$ & $++$ & $1,\ c,\ c,\ c$ \\
\end{longtable}

\noindent\textit{[The above is a representative excerpt showing the structure of Table I; the original prints all the reduced forms, characters, and composition symbols for every $D$ from $-1$ to $-100$ in three side-by-side columns spread over pp.~143--147. The continuation below abridges the per-$D$ data, listing for each $D$ the class number and a leading reduced form, to convey the bulk and structure of the table without re-typesetting every entry.]}

\begin{longtable}{r|c|p{0.7\textwidth}}
$D$ & Class no.\ & Leading reduced form(s) \\
\hline
\endhead
$-31$ & $3$ & $(1,0,31);\ (5,2,7);\ (5,-2,7)$ \\
$-32$ & $3$ & $(1,0,32);\ (4,2,9);\ (4,-2,9)$ \\
$-33$ & $4$ & $(1,0,33);\ (3,0,11);\ (2,1,17);\ (6,3,7)$ \\
$-34$ & $4$ & $(1,0,34);\ (2,0,17);\ (5,1,7);\ (5,-1,7)$ \\
$-35$ & $2$ & $(1,0,35);\ (3,1,12);\ \ldots$ \\
$-36$ & $4$ & $(1,0,36);\ (4,0,9);\ (5,2,8);\ (5,-2,8)$ \\
$-37$ & $2$ & $(1,0,37);\ (2,1,19)$ \\
$-38$ & $4$ & $(1,0,38);\ (2,0,19);\ (3,1,13);\ (3,-1,13)$ \\
$-39$ & $4$ & $(1,0,39);\ (3,0,13);\ (5,1,8);\ (5,-1,8);\ (2,1,20);\ (6,3,8)$ \\
$-40$ & $4$ & $(1,0,40);\ (4,2,11);\ (5,0,8);\ (7,2,7);\ (7,-2,7)$ \\
$-41$ & $4$ & $(1,0,41);\ (5,2,9);\ (5,-2,9);\ (2,1,21);\ (3,1,14);\ \ldots$ \\
$-42$ & $4$ & $(1,0,42);\ (2,0,21);\ (3,0,14);\ (6,0,7)$ \\
$-43$ & $2$ & $(1,0,43);\ (4,1,11);\ (4,-1,11)$ \\
$-44$ & $4$ & $(1,0,44);\ (4,0,11);\ (5,1,9);\ (5,-1,9);\ (3,1,15);\ (3,-1,15)$ \\
$-45$ & $4$ & $(1,0,45);\ (5,0,9);\ (2,1,23);\ (7,2,7)$ \\
$-46$ & $4$ & $(1,0,46);\ (2,0,23);\ (5,2,10);\ (5,-2,10)$ \\
$-47$ & $5$ & $(1,0,47);\ (3,1,16);\ (3,-1,16);\ (7,3,8);\ (7,-3,8);\ (2,1,24);\ (6,1,8);\ (4,1,12)$ \\
$-48$ & $4$ & $(1,0,48);\ (3,0,16);\ (7,1,7);\ (4,2,13);\ (4,-2,13)$ \\
$-49$ & $4$ & $(1,0,49);\ (2,1,25);\ (5,1,10);\ (5,-1,10)$ \\
$-50$ & $4$ & $(1,0,50);\ (2,0,25);\ (6,2,9);\ (6,-2,9)$ \\
$-51$ & $4$ & $(1,0,51);\ (3,1,17);\ (3,-1,17);\ (5,1,11);\ (5,-1,11);\ (4,1,13);\ (4,-1,13)$ \\
$-52$ & $4$ & $(1,0,52);\ (4,0,13);\ (3,0,17);\ (7,3,12);\ (7,-3,12);\ (2,1,26);\ (6,3,10)$ \\
$-53$ & $4$ & $(1,0,53);\ (7,2,8);\ (7,-2,8);\ (2,1,27);\ (3,1,18);\ (3,-1,18);\ (6,1,9);\ (6,-1,9)$ \\
$-54$ & $4$ & $(1,0,54);\ (2,0,27);\ (7,3,9);\ (7,-3,9);\ (5,1,11);\ (5,-1,11);\ (3,0,18);\ (6,3,11)$ \\
$-55$ & $4$ & $(1,0,55);\ (5,0,11);\ (7,1,8);\ (7,-1,8);\ (2,1,28);\ (8,3,8);\ (4,1,14);\ (4,-1,14)$ \\
$-56$ & $4$ & $(1,0,56);\ (3,1,19);\ (3,-1,19);\ (8,4,9);\ (5,2,12);\ (5,-2,12);\ (4,2,15);\ (4,-2,15);\ (2,1,28)$ \\
$-57$ & $4$ & $(1,0,57);\ (2,1,29);\ (7,0,9);\ (3,1,19);\ (3,-1,19);\ (6,3,11);\ (6,-3,11);\ (8,1,9);\ (8,-1,9)$ \\
$-58$ & $2$ & $(1,0,58);\ (2,0,29);\ (3,0,19);\ \ldots$ \\
$-59$ & $3$ & $(1,0,59);\ (3,1,20);\ (3,-1,20);\ (4,-1,15);\ (4,1,15);\ (5,1,12);\ (5,-1,12)$ \\
$-60$ & $4$ & $(1,0,60);\ (4,0,15);\ (3,0,20);\ (5,0,12);\ (7,2,10);\ (7,-2,10);\ (6,-1,10);\ (6,1,10)$ \\
$-61$ & $2$ & $(1,0,61);\ (5,2,13);\ (5,-2,13)$ \\
$-62$ & $4$ & $(1,0,62);\ (2,0,31);\ (7,3,9);\ (7,-3,9);\ (9,4,8);\ (9,-4,8);\ (6,1,11);\ (6,-1,11)$ \\
$-63$ & $4$ & $(1,0,63);\ (7,0,9);\ (8,3,9);\ (8,-3,9);\ (2,1,32);\ (6,3,12);\ \ldots$ \\
$-64$ & $4$ & $(1,0,64);\ (8,4,10);\ (5,2,13);\ (5,-2,13);\ \ldots$ \\
$-65$ & $4$ & $(1,0,65);\ (5,0,13);\ (3,1,22);\ (3,-1,22);\ (2,1,33);\ (7,1,10);\ (7,-1,10);\ (9,2,8);\ (9,-2,8)$ \\
$-66$ & $4$ & $(1,0,66);\ (2,0,33);\ (3,0,22);\ (6,0,11);\ (5,2,14);\ (5,-2,14);\ (7,2,10);\ (7,-2,10)$ \\
$-67$ & $1$ & $(1,0,67)$ \\
$-68$ & $4$ & $(1,0,68);\ (4,0,17);\ (8,2,9);\ (8,-2,9);\ (3,1,23);\ (3,-1,23);\ (7,3,11);\ (7,-3,11);\ (2,1,34)$ \\
$-69$ & $4$ & $(1,0,69);\ (3,0,23);\ (6,3,13);\ (5,1,14);\ (5,-1,14);\ (2,1,35);\ \ldots$ \\
$-70$ & $4$ & $(1,0,70);\ (2,0,35);\ (5,0,14);\ (7,0,10);\ \ldots$ \\
$-71$ & $7$ & $(1,0,71);\ (3,1,24);\ (3,-1,24);\ (5,2,15);\ (5,-2,15);\ (8,1,9);\ (8,-1,9);\ (4,1,18);\ (4,-1,18);\ (6,1,12)$ \\
$-72$ & $4$ & $(1,0,72);\ (4,2,19);\ (4,-2,19);\ (8,0,9);\ (6,0,12);\ (3,0,24)$ \\
$-73$ & $4$ & $(1,0,73);\ (8,4,11);\ (2,1,37);\ (7,2,11);\ (7,-2,11)$ \\
$-74$ & $4$ & $(1,0,74);\ (3,1,25);\ (3,-1,25);\ (9,4,10);\ (9,-4,10);\ (6,2,13);\ (6,-2,13);\ (5,1,15);\ (5,-1,15)$ \\
$-75$ & $4$ & $(1,0,75);\ (4,1,19);\ (4,-1,19);\ (7,3,12);\ (7,-3,12);\ (5,0,15);\ (3,0,25);\ (2,1,38);\ (6,3,14)$ \\
$-76$ & $4$ & $(1,0,76);\ (4,0,19);\ (5,2,16);\ (5,-2,16);\ (7,1,11);\ (7,-1,11)$ \\
$-77$ & $4$ & $(1,0,77);\ (9,2,9);\ (2,1,39);\ (3,1,26);\ (3,-1,26);\ (6,1,13);\ (6,-1,13);\ (7,0,11)$ \\
$-78$ & $4$ & $(1,0,78);\ (2,0,39);\ (3,0,26);\ (6,0,13);\ \ldots$ \\
$-79$ & $5$ & $(1,0,79);\ (3,0,26\,\text{var.});\ (5,1,16);\ (5,-1,16);\ (8,3,11);\ (8,-3,11);\ (4,1,20);\ (4,-1,20);\ (2,1,40)$ \\
$-80$ & $4$ & $(1,0,80);\ (9,1,9);\ (3,1,27);\ (3,-1,27);\ (5,0,16);\ (4,2,21);\ (4,-2,21);\ (8,1,10);\ (8,-1,10);\ (7,2,12);\ (7,-2,12)$ \\
$-81$ & $3$ & $(1,0,81);\ (9,3,10);\ (9,-3,10);\ (5,2,17);\ (5,-2,17);\ (2,1,41)$ \\
$-82$ & $4$ & $(1,0,82);\ (2,0,41);\ (7,3,13);\ (7,-3,13);\ (3,1,28);\ (3,-1,28);\ (9,4,11);\ (9,-4,11);\ (4,1,21);\ (4,-1,21);\ \ldots$ \\
$-83$ & $3$ & $(1,0,83);\ (3,1,28);\ (3,-1,28);\ (7,1,12);\ (7,-1,12);\ (4,1,21);\ (4,-1,21)$ \\
$-84$ & $4$ & $(1,0,84);\ (4,0,21);\ (3,0,28);\ (7,0,12);\ (6,0,14);\ (5,2,17);\ (5,-2,17);\ (10,4,9);\ (10,-4,9);\ (2,1,42);\ (6,1,14);\ (6,-1,14)$ \\
$-85$ & $4$ & $(1,0,85);\ (5,0,17);\ (2,1,43);\ (10,5,11);\ (7,3,13);\ (7,-3,13);\ (5,1,17);\ (5,-1,17)$ \\
$-86$ & $4$ & $(1,0,86);\ (2,0,43);\ (5,2,18);\ (5,-2,18);\ (6,2,15);\ (6,-2,15);\ (9,2,10);\ (9,-2,10);\ (3,1,29);\ (3,-1,29)$ \\
$-87$ & $6$ & $(1,0,87);\ (3,0,29);\ (4,1,22);\ (4,-1,22);\ (8,1,11);\ (8,-1,11);\ (7,2,13);\ (7,-2,13);\ (8,-3,12);\ (8,3,12);\ (2,1,44);\ (6,3,16)$ \\
$-88$ & $4$ & $(1,0,88);\ (4,2,23);\ (4,-2,23);\ (8,0,11);\ (8,4,13);\ \ldots$ \\
$-89$ & $4$ & $(1,0,89);\ (2,1,45);\ (3,1,30);\ (3,-1,30);\ (5,1,18);\ (5,-1,18);\ (9,1,10);\ (9,-1,10);\ (6,1,15);\ (6,-1,15);\ (7,3,14);\ (7,-3,14)$ \\
$-90$ & $4$ & $(1,0,90);\ (2,0,45);\ (5,0,18);\ (9,0,10);\ (7,1,13);\ (7,-1,13);\ (9,3,11);\ (9,-3,11)$ \\
$-91$ & $4$ & $(1,0,91);\ (4,1,23);\ (4,-1,23);\ (7,0,13);\ (5,2,19);\ (5,-2,19);\ (10,3,10);\ (2,1,46)$ \\
$-92$ & $4$ & $(1,0,92);\ (9,4,12);\ (9,-4,12);\ (4,0,23);\ (3,1,31);\ (3,-1,31)$ \\
$-93$ & $4$ & $(1,0,93);\ (3,0,31);\ (2,1,47);\ (6,3,17);\ (7,2,14);\ (7,-2,14);\ (5,1,19);\ (5,-1,19)$ \\
$-94$ & $4$ & $(1,0,94);\ (2,0,47);\ (5,0,19);\ (10,4,11);\ (10,-4,11);\ (9,2,11);\ (9,-2,11);\ (3,1,32);\ (3,-1,32);\ (8,3,13);\ (8,-3,13)$ \\
$-95$ & $8$ & $(1,0,95);\ (4,1,24);\ (4,-1,24);\ (5,1,12);\ (10,5,12);\ (8,1,12);\ (8,-1,12);\ (6,1,16);\ (6,-1,16)$ \\
$-96$ & $4$ & $(1,0,96);\ (4,2,25);\ (5,2,20);\ (5,-2,20);\ (3,0,32);\ (11,5,11);\ (7,3,15);\ (7,-3,15)$ \\
$-97$ & $2$ & $(1,0,97);\ (2,1,49);\ (7,1,14);\ (7,-1,14)$ \\
$-98$ & $4$ & $(1,0,98);\ (9,1,11);\ (9,-1,11);\ (2,0,49);\ \ldots$ \\
$-99$ & $4$ & $(1,0,99);\ (4,-1,25);\ (4,1,25);\ (5,1,20);\ (5,-1,20);\ (9,0,11);\ (3,1,33);\ (6,2,17);\ (6,-2,17);\ (3,-1,33)$ \\
$-100$ & $4$ & $(1,0,100);\ (4,0,25);\ (8,2,13);\ (8,-2,13);\ (10,1,10);\ (2,1,50)$ \\
\end{longtable}

\subsection*{Table II. Formes quadratiques binaires ayant pour d\'eterminants les nombres positifs non-carr\'es depuis $D = 2$ jusqu'\`a $D = 99$ \emph{(commencement; suite p.~151)}.}

\medskip

\noindent\textit{[For each positive determinant $D$, Cayley tabulates the reduced forms in their classes, with the characters $\alpha,\,\beta,\,\delta,\,\epsilon,\,\delta\epsilon$ as appropriate, the composition symbol $Cp$, and the periods of the reduced forms. The periods are printed in the abbreviated style described in the introduction: a string of integers whose successive pairs $(a_k,\,b_k,\,c_k)$, $(a_{k+1},\,b_{k+1},\,c_{k+1})$, $\ldots$ share their boundary coefficients, with $a$ and $c$ in larger type and $b$ in smaller type. The present excerpt resumes Table II from $D = 2$ to $D = 61$; the continuation $D = 62$--$99$ appears on pp.~151--156 of the next instalment.]}

\medskip

\begin{longtable}{r|c|p{0.7\textwidth}}
$D$ & Classes & Reduced forms; periods \\
\hline
\endhead
$2$ & $1$ & $(1,0,-2)$; period $1,1,-1,1,1$ \\
$3$ & $1$ & $(1,0,-3)$; period $1,1,-2,1,1$ \\
$5$ & $1$ & $(1,0,-5)$; period $1,2,-1,2,1$ \\
$6$ & $1$ & $(1,0,-6)$; period $1,2,-2,2,1$ \\
$7$ & $1$ & $(1,0,-7)$; period $1,2,-3,1,2,1,-3,2,1$ \\
$8$ & $1$ & $(1,0,-8)$; period $1,2,-4,2,1$ \\
$10$ & $2$ & $(1,0,-10);\ (2,0,-5)$; periods $1,3,-1,3,1$;\ $2,2,-3,1,3,2,-2,2,3,1,-3,2,2$ \\
$11$ & $1$ & $(1,0,-11)$; period $1,3,-2,3,1$ \\
$12$ & $1$ & $(1,0,-12)$; period $1,3,-3,3,1$ \\
$13$ & $1$ & $(1,0,-13)$; period $1,3,-4,1,3,2,-3,1,4,3,-1,3,4,1,-3,2,3,1,-4,3,1$ \\
$14$ & $1$ & $(1,0,-14)$; period $1,3,-5,2,2,2,-5,3,1$ \\
$15$ & $2$ & $(1,0,-15);\ (3,0,-5)$; periods $1,3,-6,3,1$;\ $2,3,-2,3,2$ \\
$17$ & $1$ & $(1,0,-17)$; period $1,4,-1,4,1$ \\
$18$ & $1$ & $(1,0,-18)$; period $1,4,-2,4,1$ \\
$19$ & $1$ & $(1,0,-19)$; period $1,4,-3,2,5,3,-2,3,5,2,-3,4,1$ \\
$20$ & $1$ & $(1,0,-20)$; period $1,4,-4,4,1$ \\
$21$ & $2$ & $(1,0,-21);\ (3,0,-7)$; periods $1,4,-5,1,4,3,-3,3,4,1,-5,4,1$;\ $2,3,-6,3,2$ \\
$22$ & $1$ & $(1,0,-22)$; period $1,4,-6,2,3,4,-2,4,3,2,-6,4,1$ \\
$23$ & $1$ & $(1,0,-23)$; period $1,4,-7,3,2,3,-7,4,1$ \\
$24$ & $2$ & $(1,0,-24);\ (3,0,-8)$; periods $1,4,-8,4,1$;\ $3,3,-5,2,4,2,-5,3,3$ \\
$26$ & $1$ & $(1,0,-26)$; period $1,5,-1,5,1$ \\
$27$ & $1$ & $(1,0,-27)$; period $1,5,-2,5,1$ \\
$28$ & $1$ & $(1,0,-28)$; period $1,5,-3,4,4,4,-3,5,1$ \\
$29$ & $1$ & $(1,0,-29)$; period $1,5,-4,3,5,2,-5,3,4,5,-1,5,4,3,-5,2,5,3,-4,5,1$ \\
$30$ & $2$ & $(1,0,-30);\ (2,0,-15)$; periods $1,5,-5,5,1$;\ $2,4,-7,3,3,3,-7,4,2$ \\
$31$ & $1$ & $(1,0,-31)$; period $1,5,-6,1,5,4,-3,5,2,5,-3,4,5,1,-6,5,1$ \\
$32$ & $1$ & $(1,0,-32)$; period $1,5,-7,2,4,2,-7,5,1$ \\
$33$ & $2$ & $(1,0,-33);\ (3,0,-11)$; periods $1,5,-8,3,3,3,-8,5,1$;\ $2,5,-4,3,6,3,-4,5,2$ \\
$34$ & $2$ & $(1,0,-34);\ (3,0,-\!\sim)$\ \textit{(2 ambig.\ class)}; periods $1,5,-9,4,2,4,-9,5,1$;\ $3,5,-3,4,6,2,-5,3,5,2,-6,4,3$ \\
$35$ & $1$ & $(1,0,-35)$; period $1,5,-10,5,1$ \\
$37$ & $2$ & $(1,0,-37);\ (3,1,-12)$; periods $1,6,-1,6,1$;\ $3,4,-7,3,4,5,-3,4,7,3,-4,5,3$ \\
$38$ & $1$ & $(1,0,-38)$; period $1,6,-2,6,1$ \\
$39$ & $2$ & $(1,0,-39);\ (3,0,-13)$; periods $1,6,-3,6,1$;\ $2,5,-7,2,5,3,-5,3,5,2,-7,5,2$ \\
$40$ & $2$ & $(1,0,-40);\ (3,1,-13)$; periods $1,6,-4,6,1$;\ $3,4,-8,4,3,5,-5,5,3$ \\
$41$ & $1$ & $(1,0,-41)$; period $1,6,-5,4,5,6,-1,6,5,4,-5,6,1$ \\
$42$ & $2$ & $(1,0,-42);\ (2,0,-21)$; periods $1,6,-6,6,1$;\ $2,6,-3,6,2$ \\
$43$ & $1$ & $(1,0,-43)$; period $1,6,-7,1,6,5,-3,4,9,5,-2,5,9,4,-3,5,6,1,-7,6,1$ \\
$44$ & $1$ & $(1,0,-44)$; period $1,6,-8,2,5,3,-7,4,4,4,-7,3,5,2,-8,6,1$ \\
$45$ & $2$ & $(1,0,-45);\ (2,1,-22)$; periods $1,6,-9,3,4,5,-5,5,4,3,-9,6,1$;\ $2,5,-10,5,2$ \\
$46$ & $1$ & $(1,0,-46)$; period $1,6,-10,4,3,6,-7,1,6,4,-5,6,2,5,-5,4,6,5,-7,6,3,4,-10,6,1$ \\
$47$ & $1$ & $(1,0,-47)$; period $1,6,-11,5,2,5,-11,6,1$ \\
$48$ & $2$ & $(1,0,-48);\ (3,0,-16)$; periods $1,6,-12,6,1$;\ $3,6,-4,6,3$ \\
$50$ & $1$ & $(1,0,-50)$; period $1,7,-1,7,1$ \\
$51$ & $2$ & $(1,0,-51);\ (3,0,-17)$; periods $1,7,-2,7,1$;\ $3,6,-5,4,7,3,-6,5,7,4,-5,6,3$ \\
$52$ & $1$ & $(1,0,-52)$; period $1,7,-3,6,9,4,-4,4,9,5,-3,7,1$ \\
$53$ & $1$ & $(1,0,-53)$; period $1,7,-4,6,7,2,-7,6,4,7,-1,7,4,5,-7,2,7,5,-4,7,1$ \\
$54$ & $2$ & $(1,0,-54);\ (2,1,-27)$; periods $1,7,-5,3,9,6,-2,6,9,3,-5,7,1$;\ $2,7,-2,7,2$ \\
$55$ & $2$ & $(1,0,-55);\ (3,1,-18)$; periods $1,7,-6,5,5,5,-6,7,1$;\ $2,7,-3,5,10,3,-3,7,2$ \\
$56$ & $2$ & $(1,0,-56);\ (3,2,-13)$; periods $1,7,-7,7,1$;\ $4,6,-5,4,8,4,-5,6,4$ \\
$57$ & $2$ & $(1,0,-57);\ (3,0,-19)$; periods $1,7,-8,1,7,5,-3,5,7,1,-8,7,1$;\ $2,7,-4,5,8,3,-6,3,8,5,-4,7,2$ \\
$58$ & $1$ & $(1,0,-58)$; period $1,7,-9,2,6,4,-7,3,7,4,-6,2,9,7,-1,7,9,2,-6,4,7,3,-7,4,6,2,-9,7,1$ \\
$59$ & $1$ & $(1,0,-59)$; period $1,7,-10,3,5,7,-2,7,5,3,-10,7,1$ \\
$60$ & $2$ & $(1,0,-60);\ (3,0,-20)$; periods $1,7,-11,4,4,4,-11,7,1$;\ $3,6,-8,2,7,6,-5,6,7,2,-8,6,3$ \\
$61$ & $1$ & $(1,0,-61)$; period $1,7,-12,5,3,7,-4,6,9,4,-5,8,5,4,-9,6,4,7,-3,5,12,7,1$ \\
\end{longtable}

\noindent\textit{[Table II continues to $D = 99$ on pp.~151--156 of the next instalment; Table III, the thirteen irregular negative determinants of the first thousand, follows there as well. The structure of the entries is preserved throughout: classes, characters, composition symbol $Cp$, and periods in the abbreviated notation defined above.]}

\medskip
\noindent\textit{Londres, 6 Novembre, 1860.}

\medskip
\noindent[Continues on p.~151 with Table II, $D = 62$ to $99$, and Table III.]

\end{document}
